% Origen: Secciones/PI - Conclusiones y Futuros Pasos.md
\chapter{Conclusiones y futuros pasos}
\label{C:conclusiones}

% Figuras de ESTE capítulo + las compartidas. Siempre al principio del archivo.
\graphicspath{{figs/cap11/}{figs/comunes/}}

%==============================================================================
\section{Conclusiones}
\label{S:concl}

Implementamos un analizador multicanal y un osciloscopio por eventos sobre la
Red Pitaya STEMlab 125-14. Todo el procesamiento de los pulsos corre en la PL:
el disparo, la medición de la amplitud y de la forma, el histograma y la captura
de las trazas. El PS se queda con el control, la adquisición y el guardado. El
sistema final es el TOP de integración con el anillo de eventos
(Cap.~\ref{C:top}), y lo caracterizamos
con \codigo{gen\_pulsos} en el lazo digital y en el canal analógico
OUT1 $\rightarrow$ IN1 (Cap.~\ref{C:resultados}).

\subsection{Parámetros del sistema implementado}
\label{S:concl-parametros}

En la Tabla~\ref{t:concl-parametros} se presentan los parámetros fijados en la
síntesis, que sólo cambian volviendo a sintetizar, y la ocupación de la FPGA.

\begin{table}[H]
\centering
\caption[Parámetros del sistema implementado]{%
Parámetros del sistema implementado: los que fija la síntesis del TOP de
integración con el anillo de eventos y su ocupación de la
Zynq-7010.\label{t:concl-parametros}}
\small
\begin{tabularx}{\textwidth}{|>{\raggedright\arraybackslash}p{0.16\textwidth}%
  |>{\raggedright\arraybackslash}p{0.22\textwidth}%
  |>{\raggedright\arraybackslash}X|}
\hline \encabezado \textbf{Bloque} & \textbf{Parámetro} & \textbf{Valor} \\ \hline
Plataforma & placa & Red Pitaya STEMlab 125-14: Zynq-7010, con dos núcleos ARM
  Cortex-A9 \\ \hline
 & ADC y DAC & 14 bits, 125~MS/s (\qty{8}{\nano\second} por muestra); entrada
  en LV, $\pm\qty{1}{\volt}$ (1~ADU $\approx$ \qty{122}{\micro\volt}) \\ \hline
MCA & espectro 1D & 8192 canales, con contadores de 32 bits que saturan
  \\ \hline
 & mapa 2D & $128 \times 64$ canales (amplitud $\times$ forma) \\ \hline
 & ranuras de medición & 3 \\ \hline
 & banco de trapecios & 3 formas, $(k, m) = (64, 48)$, $(32, 24)$ y $(8, 8)$,
  con un polo-cero común \\ \hline
 & publicador & uno, de 18 ciclos por evento \\ \hline
 & contabilidad del tiempo & \codigo{realtime} = \codigo{vetotime} +
  \codigo{t\_win} + \codigo{t\_div} + \codigo{t\_rearm} + \codigo{livetime}, al
  ciclo \\ \hline
Multitrigger & fuentes de disparo & nivel, pendiente ($L = 8$, $G = 4$),
  polo-cero ($k = 2$, $m = 1$), externa, generador y software, combinadas con
  una máscara OR \\ \hline
Osciloscopio & captura & por la BRAM, un evento por vez, y por el anillo de
  eventos en la DDR \\ \hline
 & ventanas del anillo & de $8 + 24$ a $32 + 992$ muestras (antes + después
  del disparo) \\ \hline
Generador & \codigo{gen\_pulsos} & tabla de 1024 alturas; caída, subida y ruido
  blanco programables; arribos periódicos o de Poisson; lazo digital \\ \hline
FPGA & LUT & 14\,358 de 17\,600 (\qty{81.6}{\percent}) \\ \hline
 & registros & 15\,513 de 35\,200 (\qty{44.1}{\percent}) \\ \hline
 & slices & 4400 de 4400 (\qty{100}{\percent}) \\ \hline
 & BRAM (bloques de 36~kb) & 33 de 60 (\qty{55}{\percent}) \\ \hline
 & DSP & 10 de 80 (\qty{12.5}{\percent}) \\ \hline
 & timing & cierra a \qty{125}{\mega\hertz}, con WNS = \qty{+0.008}{\nano\second}
  y WHS = \qty{+0.040}{\nano\second} \\ \hline
\end{tabularx}
\end{table}

El diseño ocupa las 4400 slices de la Zynq-7010, con el \qty{81.6}{\percent} de
las LUT, y cierra el timing a \qty{125}{\mega\hertz} con
\qty{8}{\pico\second} de margen. Para hacer lugar al anillo y a
\codigo{gen\_pulsos} sacamos la captura del osciloscopio por AXI, que el anillo
reemplaza, y bajamos el MCA de 4 a 3 ranuras y de 4 a 3 formas, lo que costó
cerca del \qty{1}{\percent} de la tasa en la línea. La placa quedó llena: un
bloque nuevo pide sacar otro o pasar a una FPGA más grande.

\subsection{Desempeño del hardware}
\label{S:concl-hw}

En la Tabla~\ref{t:concl-hw} se resumen las figuras de mérito que medimos con
\codigo{gen\_pulsos}. Salvo que se diga otra cosa, son con 3 ranuras, el banco
de trapecios y la base fija.

\begin{table}[H]
\centering
\caption[Figuras de mérito del MCA medidas con \codigo{gen\_pulsos}]{%
Figuras de mérito del MCA medidas con \codigo{gen\_pulsos}, en el lazo digital
y en el canal analógico OUT1 $\rightarrow$ IN1 con \qty{50}{\ohm}. Las
amplitudes van en ADU y las posiciones del histograma en canales. ---: todavía
sin medir en ese modo.\label{t:concl-hw}}
\small
\begin{tabularx}{\textwidth}{|>{\raggedright\arraybackslash}p{0.27\textwidth}%
  |>{\raggedright\arraybackslash}X|>{\raggedright\arraybackslash}X|}
\hline \encabezado \textbf{Figura de mérito} & \textbf{Lazo digital} &
  \textbf{Canal analógico} \\ \hline
Tasa aceptada, con 3~M/s de entrada & 1.9~M/s & casi la misma \\ \hline
Tasa en la línea, con 3~M/s de entrada & 1.2~M/s con el ASG y el polo-cero;
  0.78~M/s con la pendiente & máximo de 0.43 a 0.53~M/s, con 1.5~M/s de entrada
  \\ \hline
Error de $\hat\lambda$, con 3~M/s de entrada & ASG $-\qty{1.7}{\percent}$;
  polo-cero $-\qty{3.7}{\percent}$; pendiente $-\qty{18}{\percent}$ &
  polo-cero $+\qty{5.1}{\percent}$, por los disparos de más sobre el sobrepico
  \\ \hline
Separación de los disparos de un par & polo-cero \qty{40}{\nano\second};
  pendiente \qty{160}{\nano\second} & pendiente igual; el polo-cero vuelve a
  disparar \\ \hline
Separación de un par con las dos amplitudes en su línea &
  \qty{128}{\nano\second} (polo-cero) y \qty{160}{\nano\second} (pendiente);
  \qty{1.5}{\micro\second} con 1 ranura & pendiente igual; el polo-cero no lo
  resuelve \\ \hline
Umbral (mitad de los pulsos disparan) & polo-cero 48~ADU; pendiente 81~ADU &
  51 y 82~ADU \\ \hline
Walk del disparo & polo-cero \qty{8}{\nano\second}; pendiente
  \qty{48}{\nano\second} & 16 y \qty{55}{\nano\second} \\ \hline
INL del trapecio & $<\qty{0.002}{\percent}$~FS & $\pm\qty{0.008}{\percent}$~FS
  hasta 3200~ADU; $+\qty{0.66}{\percent}$ en 8000, donde satura el ADC
  \\ \hline
DNL del MCA & $<\qty{1.5}{\percent}$ rms (cota: domina la tabla del generador)
  & --- \\ \hline
Ganancia de ruido, $k = 8$, 32 y 64 & 0.48, 0.24 y 0.18 & igual \\ \hline
Piso de ruido $\sigma_0$, $k = 64$ & 0.06~ADU & 1.41~ADU \\ \hline
Resolución en una línea de 1500~ADU, $k = 64$ y ruido 2.3~ADU & FWHM
  \qty{0.07}{\percent} & FWHM \qty{0.24}{\percent} \\ \hline
Corrimiento de la línea con la tasa & $< 1.5$ canales hasta 3~M/s &
  $\sim 9$ canales a 400~k/s \\ \hline
Ajuste del polo-cero & error nulo en $\Delta M = +0.03$ a
  $+\qty{0.06}{\percent}$ & --- \\ \hline
\end{tabularx}
\end{table}

De estas mediciones sacamos cinco conclusiones.

\textbf{La cadena digital no agrega un error medible.} En el lazo digital la INL
del trapecio queda por debajo de \qty{0.002}{\percent}~FS y las líneas se
corren menos de 1.5 canales hasta 3~M/s. Lo que se aparta del ideal en el canal
analógico es del canal, no del diseño. Esto también resuelve una limitación de
la generación externa controlada, que no separaba la INL del generador de la del
ADC: el canal de la placa da $\pm\qty{0.008}{\percent}$~FS hasta 3200~ADU, así
que lo que mida de más un generador externo en ese rango es del generador.

\textbf{Lo que limita la tasa es el apilamiento, no el tiempo muerto.} A 3~M/s
de entrada, pasar de 1 a 3 ranuras sube los eventos aceptados de 0.73 a
1.9~M/s: las ranuras sacan del camino casi todo el tiempo muerto. Lo que queda
es la diferencia entre los aceptados y los que caen en su línea, 1.9 contra
1.2~M/s, y es casi toda apilamiento: pulsos que el MCA mide con otro encima. Esa
pérdida es estadística, depende de la tasa y del largo de la forma, y no baja
con más ranuras.

\textbf{El disparo fija el comportamiento a alta tasa.} El polo-cero separa dos
disparos a \qty{40}{\nano\second} y registra en su línea lo mismo que el ASG,
que dispara con cada pulso sin mirar la señal (1.23 contra 1.21~M/s a 3~M/s). La
pendiente necesita \qty{160}{\nano\second} y es paralizable: su tasa en la línea
llega a 0.97~M/s con 2~M/s de entrada y baja a 0.78 con 3~M/s. Eso llega a la
tasa estimada: $\hat\lambda$ corrige el tiempo muerto de las ranuras y del
publicador, y queda a $-\qty{1.7}{\percent}$ de la tasa real con el ASG y a
$-\qty{3.7}{\percent}$ con el polo-cero, pero con la pendiente queda
\qty{18}{\percent} abajo, porque el re-armado del disparo sólo se cuenta cuando
hay una ranura libre. Con un canal sin sobrepico, el polo-cero es la fuente de
disparo para medir.

\textbf{En el canal analógico, lo que limita es el canal.} El canal
OUT1 $\rightarrow$ IN1 con \qty{50}{\ohm} tiene una ganancia de 0.95, cae
\qty{3}{\decibel} en \qty{52}{\mega\hertz} y tiene un piso de 4.7~ADU rms. Dos
rasgos suyos explican casi todo lo que el modo analógico agrega:
\begin{itemize}
  \item \textbf{El sobrepico del \qty{24}{\percent}.} El polo-cero vuelve a
    disparar sobre él desde unas 3200~ADU y deja de resolver pares, y desde unas
    6990~ADU el ADC satura. La pendiente no lo ve.
  \item \textbf{La caída lenta del \qty{7.3}{\percent}, en 3 a
    \qty{13}{\micro\second}.} Las colas negativas de los pulsos se apilan y
    bajan la base con la tasa, y el trapecio pasa ese error a la amplitud. Con
    base fija la línea se corre unos 9 canales a 400~k/s, y la tasa en la línea
    no pasa de 0.43 a 0.53~M/s.
\end{itemize}

Estos dos rasgos son del montaje de generación interna, y todavía no separamos
cuánto viene de la salida y cuánto de la entrada; con un detector sólo se usa la
entrada. La conclusión, igual, vale para cualquier canal: antes de medir hay que
caracterizarlo, y una cola lenta hay que compensarla antes del MCA, como el
polo-cero compensa la caída del pulso.

\subsection{Desempeño del software}
\label{S:concl-sw}

En el software medimos cuánto cuesta cada capa del camino de los datos. Los
resultados se resumen en la Tabla~\ref{t:concl-sw}.

\begin{table}[H]
\centering
\caption[Desempeño del software en la placa]{%
Desempeño del software en la placa, por capa del camino de los datos. «Al día»
es la tasa más alta del barrido sin descartes ni atraso; «saturado», lo que lee
la capa cuando no da abasto.\label{t:concl-sw}}
\small
\begin{tabularx}{\textwidth}{|>{\raggedright\arraybackslash}p{0.33\textwidth}%
  |>{\raggedright\arraybackslash}X|}
\hline \encabezado \textbf{Camino} & \textbf{Medido} \\ \hline
leer un registro & \qty{1.5}{\micro\second} en C, \qty{56.8}{\micro\second} por
  el binding y \qty{10.5}{\micro\second} en Python puro \\ \hline
leer el espectro (8192 canales) & \qty{1.70}{\milli\second} en C,
  \qty{1.88}{\milli\second} por el binding y \qty{34.3}{\milli\second} en Python
  puro; \qty{3.5}{\milli\second} por la red \\ \hline
osciloscopio por la BRAM, con el lazo en C & \qty{18.2}{\micro\second} de tiempo
  muerto por evento (mediana); 34~k/s guardados \\ \hline
anillo, la PL sola & 1.42~M/s con ventanas de $8 + 24$ muestras; 49~k/s con
  $32 + 992$ \\ \hline
anillo, el lector de C copiando & 753~k/s con $8 + 24$ \\ \hline
anillo, el lector de C copiando y verificando & 300~k/s al día con $8 + 24$
  (439~k/s saturado); de 70 a 98~MB/s saturado \\ \hline
anillo, la sesión de Python guardando en \ruta{.npz} & 100~k/s al día con
  $8 + 24$ (167~k/s saturado en la SD y 203~k/s en \ruta{/dev/shm}); 10~k/s con
  $32 + 992$ \\ \hline
$10^8$ eventos con los dos núcleos ocupados & la PL sola a 1.4~M/s y el lector
  de C a 200~k/s, sin descartes ni errores \\ \hline
\end{tabularx}
\end{table}

\textbf{Lo que limita al osciloscopio es el PS.} Cada capa que se suma al camino
de los datos baja la tasa: la PL sola escribe 1.42~M/s en el anillo, el lector
de C copia 753~k/s y, verificando cada slot, 300~k/s, y la sesión de Python
guarda 100~k/s. El escritor de la PL llega en la placa al techo de la simulación
con todas las ventanas, así que el AXI y la DDR no lo frenan; con la sesión,
\ruta{/dev/shm} da casi lo mismo que la SD, así que tampoco frena el disco. Lo
que limita es el software en los dos núcleos del ARM. El MCA no tiene este
problema, porque acumula el histograma en la PL y el PS sólo lee el espectro:
acepta 1.9~M/s, contra los 100~k/s de trazas que guarda el osciloscopio.

\textbf{El costo del binding es por llamada, no por dato.} Pasar por
\codigo{ctypes} agrega unos \qty{55}{\micro\second} por llamada, que dominan al
leer un registro y casi no pesan al leer el espectro (un \qty{10}{\percent} sobre
C). Por eso los lazos por evento están en C, y en Python quedan la sesión, los
archivos y la aplicación. Con el lazo de la BRAM en C, el tiempo muerto por
evento bajó de 252.7 a \qty{18.2}{\micro\second}.

\textbf{El anillo absorbe las demoras del PS, hasta un punto.} Con $10^8$
eventos y los dos núcleos ocupados por otros procesos, la PL a 1.4~M/s no
descarta ninguno: el anillo absorbió una demora de \qty{30}{\milli\second} del
lector. El lector de C cumple a 200~k/s, pero sin margen: una demora de medio
segundo llenó el anillo al \qty{81}{\percent}. Con el ARM ocupado, la tasa
segura para leer trazas de 32 muestras desde C está por debajo de 200~k/s.

En resumen, el procesamiento en la PL rinde lo que permite la estadística de los
arribos. Lo que queda por mejorar está en el canal analógico, en la estabilidad
con la temperatura, que todavía no medimos, y en el camino de las trazas por el
PS.

%==============================================================================
\section{Posibles futuros pasos}
\label{S:concl-futuros}

A continuación se describen algunos de los pasos que quedan para llevar el
sistema a una medición real y para mejorarlo.
\begin{itemize}
  \item \textbf{Implementar el sistema en un setup experimental real.} Es la
    última etapa de validación: un detector con su preamplificador en IN1 y
    fuentes de calibración. Con pulsos reales se miden la resolución en energía,
    la linealidad y la tasa, y se ve cuánto de lo que agrega el canal analógico
    aparece también con el detector, que sólo usa la entrada.
  \item \textbf{Comparar el funcionamiento con otros MCA comerciales.} Con el
    mismo detector y la misma fuente, contra un MCA digital comercial. Antes se
    puede comparar en papel con sus especificaciones, cuidando que las cifras
    sean comparables: la INL, por ejemplo, la dan sobre casi todo el rango, y la
    nuestra está medida hasta 3200~ADU.
  \item \textbf{Obtener feedback de usuarios y comprobar la facilidad de uso de
    la API y la GUI.} Darle el sistema a usuarios del laboratorio con una tarea
    fija, como calibrar el canal, adquirir un espectro y guardar las trazas, y
    medir cuánto tardan y dónde se traban, con la GUI y con la API.
  \item \textbf{Crear una aplicación web al estilo de las de Red Pitaya.} Como un
    cliente más del protocolo JSON del servidor, el mismo que usa la GUI, sin
    cambiar nada en la placa.
  \item \textbf{Crear una documentación oficial para el código de la API.} La
    documentación ya está en el repositorio; falta publicarla como un sitio
    generado a partir del código.
  \item \textbf{Compensar la caída lenta del canal.} La base dinámica, que
    corrige la base según la tasa estimada, ya está implementada en el software y
    falta validarla en la placa. Corrige la base en promedio, pero no su
    fluctuación de un pulso a otro. Para eso hay que invertir los polos lentos
    del canal antes del MCA, con un filtro como el polo-cero del trapecio.
  \item \textbf{Estabilizar el espectro.} Primero hay que medir la deriva del
    sistema con la temperatura. Con eso, un lazo en el PS puede mantener un pico
    de referencia en su canal, corrigiendo una ganancia fina y un desplazamiento
    antes del histograma, como hacen los MCA comerciales. Para estabilizar la
    electrónica alcanza con un pulser como \codigo{gen\_pulsos}; para
    estabilizar también el detector hace falta una línea de la fuente.
  \item \textbf{Contar el re-armado del disparo en $\hat\lambda$ aunque no haya
    ranura libre.} Hoy \codigo{t\_rearm} sólo cuenta cuando hay una ranura
    libre, y con la pendiente $\hat\lambda$ queda \qty{18}{\percent} abajo a
    3~M/s. Un contador del re-armado que corra también con las ranuras ocupadas
    debería corregir el tiempo muerto del disparo entero.
  \item \textbf{Testear en otras plataformas con mayor cantidad de recursos.} La
    Zynq-7010 quedó llena. Probar el diseño en una plataforma con más recursos,
    como una con la Zynq-7020, que tiene unas tres veces más lógica, permitiría
    recuperar la ranura y la forma del banco que sacamos, y agregar los bloques
    de esta lista.
  \item \textbf{Medir la salida de las trazas por la red.} El camino del anillo
    de eventos por la red está implementado y probado sin la placa. Falta medir
    en la placa cuánto da, contra los 100~k/s de la sesión que guarda en la SD.
  \item \textbf{Agregar coincidencias y anticoincidencias.} Con dos detectores
    en IN1 e IN2, o con un detector y un blindaje activo. El análisis y el
    diseño están hechos; se llevan al RTL cuando haya un montaje que los use.
\end{itemize}

%%% Local Variables:
%%% mode: latex
%%% TeX-master: "template"
%%% End:
